% ============================================================
% 计算质子中的胶子部分子分布函数 —— 理论基础与计算工作流
% ============================================================
\documentclass[10pt,aspectratio=169]{beamer}
\usepackage{ctex}
\setCJKmainfont{Songti SC}[BoldFont=Heiti SC]
\setCJKsansfont{Heiti SC}
\setCJKmonofont{STFangsong}

\usetheme{default}
\usefonttheme{professionalfonts}
\setbeamertemplate{navigation symbols}{}
\setbeamertemplate{footline}[frame number]
\setbeamercolor{title}{fg=black}
\setbeamercolor{frametitle}{fg=black}
\setbeamercolor{structure}{fg=black}

\usepackage{amsmath,amssymb}
\usepackage{physics}
\usepackage{braket}

\newcommand{\msbar}{\overline{\mathrm{MS}}}
\newcommand{\lqcd}{\Lambda_{\mathrm{QCD}}}
\newcommand{\alphas}{\alpha_s}
\newcommand{\LaMET}{\mathrm{LaMET}}
\newcommand{\GeV}{\,\mathrm{GeV}}
\newcommand{\MeV}{\,\mathrm{MeV}}
\newcommand{\fm}{\,\mathrm{fm}}

\begin{document}

% ============================================================
% Slide 1: 物理动机——为什么需要计算胶子PDF
% ============================================================
\begin{frame}{物理动机——为什么需要从第一性原理计算胶子PDF}

  胶子PDF $g(x,\mu^2)$ 描述了在质子中找到一个携带纵向动量分数 $x$ 的胶子的概率密度。
  理解胶子分布是当代 QCD 研究的核心课题之一，原因有三：

  \vspace{4pt}
  \begin{enumerate}
    \item \textbf{质子质量起源。} 质子质量 $m_p \approx 938\MeV$，其中组成夸克的质量贡献仅约占 $1\%$（$\sim 9\MeV$）；剩余约 $99\%$ 来自胶子场的量子涨落和夸克-胶子相互作用的能量。理解胶子分布是理解质量生成机制的关键。

    \item \textbf{LHC 物理的核心输入。} Higgs 玻色子的主要产生机制——胶子-胶子融合 $gg \to H$——直接依赖于 $g(x,\mu^2)$。此外，顶夸克对产生、jets 产生等关键过程均由胶子初态主导。PDF 的不确定性是高能物理许多新物理搜索中的主导理论误差来源。

    \item \textbf{质子自旋危机。} 夸克自旋仅贡献质子总自旋的约 $30\%$，胶子螺旋度分布 $\Delta g(x,\mu^2)$ 的贡献是解决这一谜题的关键。$\Delta g$ 的理论基础与非极化 $g(x,\mu^2)$ 共享相同的算符结构。
  \end{enumerate}

  实验上通过全局 QCD 分析提取胶子 PDF，但在大 $x$ 和小 $x$ 区域仍有较大不确定性。因此，从 QCD 第一性原理出发进行\textbf{非微扰计算}具有重要价值。
\end{frame}

% ============================================================
% Slide 2: 理论基础——QCD因子化与PDF的形式定义
% ============================================================
\begin{frame}{理论基础——QCD因子化与PDF的算符定义}

  前述物理动机引出一个基本理论问题：\textbf{PDF 在 QCD 中是如何严格定义的？}

  \textbf{因子化定理 (Collins--Soper--Sterman, 1989)} 给出了答案。对于高能散射 $A+B \to C+X$：
  {\footnotesize
  \begin{equation*}
    \sigma_{AB \to C}(P_A,P_B) = \sum_{a,b=q,\bar{q},g}
    \int_0^1\! dx_a dx_b\;
    f_{a/A}(x_a,\mu_F^2)\,
    f_{b/B}(x_b,\mu_F^2)\,
    \hat{\sigma}_{ab \to C}\!\bigl(x_a P_A, x_b P_B, \alpha_s(\mu_R^2), Q^2/\mu_F^2\bigr)
  \end{equation*}}

  因子化定理的物理内涵：
  \begin{itemize}
    \item 将高能散射截面分离为\textbf{普适的非微扰 PDF} $f_{a/A}$ 与\textbf{过程依赖的微扰硬散射截面} $\hat{\sigma}_{ab}$。
    \item 因子化标度 $\mu_F$ 将 QCD 辐射分为：$k_\perp < \mu_F$（软，吸收进 PDF）和 $k_\perp > \mu_F$（硬，归入 $\hat{\sigma}$）。
    \item PDF 的\textbf{普适性}意味着同一组 PDF 可用于预言不同散射过程——这是其强大之所在。
    \item 证明基于 SCET 中 Wilson 线算符乘积展开的因子化性质，依赖于红外发散的普适可因子化性。
  \end{itemize}

  PDF 不是唯象拟合的参数——它们具有 QCD 中严格的\textbf{算符定义}，这将是下一节的内容。\vspace{-4pt}
\end{frame}

% ============================================================
% Slide 3: 胶子PDF的严格算符定义——光锥关联函数
% ============================================================
\begin{frame}{胶子PDF的光锥算符定义}

  因子化定理确立了 PDF 的理论地位。其\textbf{算符定义}由光锥关联函数给出：

  引入光锥坐标 $v^\pm = (v^0 \pm v^3)/\sqrt{2}$，类光矢量 $n^\mu = (1,0,0,-1)/\sqrt{2}$（$n^2=0$）：

  \begin{equation}
    \boxed{
      g(x, \mu^2) = \frac{1}{2\pi x P^+}
      \int_{-\infty}^{\infty} \frac{d\xi^{-}}{2\pi}\;
      e^{-i x P^+ \xi^{-}}
      \big\langle P \big|
      F_a^{+\mu}(\xi^{-} n)\;
      \mathcal{W}[\xi^{-} n, 0]\;
      F_{\mu}^{a\,+}(0)
      \big| P \big\rangle_{\xi^+ = 0,\;\boldsymbol{\xi}_\perp = 0}
    }
  \end{equation}

  各要素的物理意义：
  \begin{itemize}
    \item $F_a^{\mu\nu} = \partial^\mu A_a^\nu - \partial^\nu A_a^\mu + g f_{abc} A_b^\mu A_c^\nu$ —— 胶子场强张量，携带动量信息。
    \item $\mathcal{W}[\xi^{-}, 0] = \mathcal{P}\exp[-ig\int_0^{\xi^{-}} d\lambda\, n\!\cdot\! A(\lambda n)]$ —— 沿光锥路径的 Wilson 线，\textbf{保证规范不变性}。在光锥规范 $A^+=0$ 下退化为单位矩阵。
    \item $\mu$ —— $\msbar$ 重整化标度，减除算符定义中固有的紫外发散。
    \item 关联函数定义在类光间隔 $\xi^2 = 0$ 上，这是 PDF 具有概率诠释的几何根源。
  \end{itemize}

  至此，我们有了严格的理论定义。但一个自然的追问是：$g(x,\mu^2)$ 如何随标度 $\mu$ 变化？这引向 DGLAP 演化。
\end{frame}

% ============================================================
% Slide 4: 标度演化——DGLAP方程与劈裂函数
% ============================================================
\begin{frame}{标度演化——DGLAP方程与劈裂函数}

  PDF 的紫外重整化引入了对 $\mu$ 的依赖，其演化由 \textbf{DGLAP 方程}控制：

  \begin{equation}
    \boxed{
      \mu^2 \frac{\partial}{\partial \mu^2}
      \begin{pmatrix} q_i(x,\mu^2) \\[2pt] g(x,\mu^2) \end{pmatrix}
      = \frac{\alpha_s(\mu^2)}{2\pi}
      \sum_{j} \int_x^1 \frac{dz}{z}
      \begin{pmatrix}
        P_{q_i q_j}(z) & P_{q_i g}(z) \\[4pt]
        P_{g q_j}(z)  & P_{g g}(z)
      \end{pmatrix}
      \begin{pmatrix} q_j(x/z,\mu^2) \\[2pt] g(x/z,\mu^2) \end{pmatrix}
    }
  \end{equation}

  劈裂函数 $P_{ab}(z)$ 的物理图像：部分子 $b$ 辐射出部分子 $a$，自身动量分数降为 $z$。
  微扰展开：$P_{ab} = P_{ab}^{(0)} + \frac{\alpha_s}{2\pi}P_{ab}^{(1)} + \bigl(\frac{\alpha_s}{2\pi}\bigr)^2 P_{ab}^{(2)} + \mathcal{O}(\alpha_s^3)$，
  已知到 \textbf{NNLO} (Moch--Vermaseren--Vogt, 2004)。

  \textbf{领头阶解析表达式}（体现 QCD 色结构的简洁性）：
  {\small
  \begin{align*}
    P_{qq}^{(0)}(z) &= C_F\Bigl[\frac{1+z^2}{(1-z)_+}+\frac{3}{2}\delta(1-z)\Bigr], &
    P_{gq}^{(0)}(z) &= C_F\frac{1+(1-z)^2}{z} \\
    P_{gg}^{(0)}(z) &= 2C_A\Bigl[\frac{z}{(1-z)_+}+\frac{1-z}{z}+z(1-z)\Bigr]
                       +\frac{11C_A-4n_f T_R}{6}\delta(1-z), &
    P_{qg}^{(0)}(z) &= T_R[z^2+(1-z)^2]
  \end{align*}}
  其中 $C_A=3$, $C_F=4/3$, $T_R=1/2$。``$+$''分布定义为
  {\small $\int_0^1 dz\,[f(z)]_+ g(z) = \int_0^1 dz\,f(z)[g(z)-g(1)]$}。

  DGLAP 方程告诉我们：给定一个标度的 PDF 初值，即可预演所有标度的 PDF。但\textbf{初值从哪里来？}这需要第一性原理的非微扰计算——即格点 QCD。
\end{frame}

% ============================================================
% Slide 5: 第一性原理计算——格点QCD的基本框架
% ============================================================
\begin{frame}{第一性原理计算——格点QCD的基本框架}

  格点 QCD (Wilson, 1974) 是唯一可以从 QCD 拉氏量出发进行非微扰计算的系统性方法。

  \textbf{基本设定：} 四维 Euclidean 超立方格子，格距 $a$。Wick 旋转 $t \to -i\tau$ 后，
  QCD 配分函数成为统计力学中可 Monte Carlo 模拟的形式：
  \begin{equation*}
    Z = \int \mathcal{D}U\, \mathcal{D}\bar{\psi}\mathcal{D}\psi\;
    \exp\!\bigl[-S_E[U,\bar{\psi},\psi]\bigr]
  \end{equation*}

  Wilson 规范作用量（plaquette 形式）：
  \begin{equation*}
    S_G = \beta \sum_x \sum_{\mu<\nu} \bigl(1 - \tfrac{1}{3}\mathrm{Re}\,\mathrm{Tr}\,U_{\mu\nu}(x)\bigr),
    \quad \beta = 6/g^2
  \end{equation*}
  加上 Wilson 费米子作用量 $S_F = \sum_x \bar{\psi}(x)(\not{D}_W + m)\psi(x)$ 构成完整 QCD 作用量。

  通过 \textbf{Hybrid Monte Carlo (HMC)} 算法生成规范场组态系综 $\{U^{(n)}\}_{n=1}^{N_{\mathrm{cfg}}}$。
  物理观测量由路径积分平均得到：$\langle O \rangle \approx \frac{1}{N_{\mathrm{cfg}}} \sum_n O[U^{(n)}]$。

  \textbf{可控的系统学误差：} 格点 QCD 的精度受三个可外推的系统学限制：
  \begin{itemize}
    \item 格距外推：在多组 $\beta$ 值上模拟，取 $a \to 0$ 极限。
    \item 体积外推：在多个 $L^3 \times T$ 体积上模拟，取 $V \to \infty$ 极限。
    \item 手征外推：在多个 $m_\pi$ 值上模拟，取 $m_\pi \to m_\pi^{\mathrm{phys}} \approx 140\MeV$ 极限。
  \end{itemize}

  格点 QCD 在计算强子谱、衰变常数等方面取得了巨大成功。但将其用于 PDF 计算时，遇到了一个根本性的障碍。\vspace{-6pt}
\end{frame}

% ============================================================
% Slide 6: 核心困难——Euclidean格点与Minkowski光锥的冲突
% ============================================================
\begin{frame}{核心困难——Euclidean格点与Minkowski光锥的冲突}

  回顾第 3 页的胶子 PDF 算符定义，其关键特征是关联函数定义在\textbf{类光间隔} $\xi^2 = 0$ 上。
  然而——

  \begin{center}
  \begin{tabular}{p{0.42\textwidth}|p{0.42\textwidth}}
    \hline
    \textbf{胶子PDF所需 (Minkowski)} & \textbf{格点QCD所能 (Euclidean)} \\
    \hline
    类光间隔 $\xi^2 = 0$，$\xi = (0,\xi^-,\mathbf{0}_\perp)$ &
    类空间隔 $z^2 > 0$，$z = (0,0,0,z_3)$ \\
    实时动力学，沿光锥方向积分 &
    虚时模拟，等时关联函数 \\
    紫外发散在 $\msbar$ 减除（微扰） &
    空间分离算符含线性发散 $\sim e^{-\delta m|z|/a}$（非微扰） \\
    概率诠释成立 &
    无法直接赋予概率诠释 \\
    \hline
  \end{tabular}
  \end{center}

  \textbf{本质矛盾：} QCD 中 PDF 是 Minkowski 时空的光锥物理量，而格点 QCD 的计算工具局限于 Euclidean 时空。这是质子结构格点计算面临的最基本挑战。

  \vspace{4pt}
  这一矛盾驱动了过去十年中\textbf{两类互补的解决方案}的提出，它们构成了当前格点 PDF 计算的主流方法论框架。二者的共同思想是：用 Euclidean 格点上\textbf{可计算}的关联函数，通过理论可控的因子化关系，间接提取光锥 PDF 信息。
\end{frame}

% ============================================================
% Slide 7: 解决策略——LaMET与赝PDF的双路径框架
% ============================================================
\begin{frame}{解决策略——LaMET与赝PDF的双路径框架}

  如前所述，格点可以直接计算的是\textbf{等时空间分离}的矩阵元：
  \begin{equation*}
    h_g(z, P_z) = \big\langle P \big| F_a^{3\mu}(z)\, \mathcal{W}[z,0]\, F_{\mu}^{a\,3}(0) \big| P \big\rangle
  \end{equation*}
  其中 $\mathcal{W}[z,0] = \mathcal{P}\exp[-ig\int_0^z dz' A_z(z')]$ 是沿 $z$ 轴的空间 Wilson 线。

  \textbf{两类方法从同一个 $h_g(z, P_z)$ 出发，但采用不同的极限和数学处理：}

  \begin{enumerate}
    \item \textbf{大动量有效理论 ($\LaMET$) / 准PDF方法 (Ji, 2013)}

    先将 $h_g(z, P_z)$ 对 $z$ 做 Fourier 变换得准PDF $\tilde{g}(x, P_z)$，取 $P_z \gg \lqcd$，
    通过匹配系数 $C_{gg}$ 在 $\msbar$ 方案下映射到光锥 PDF：
    {\small
    \begin{equation*}
      g(x, \mu) = \int_x^1 \frac{dy}{|y|}\; C_{gg}\!\left(\frac{x}{y}, \frac{\mu}{P_z}\right)\;
      \tilde{g}_R(y, P_z) + \mathcal{O}\!\left(\frac{\lqcd^2}{x^2 P_z^2}\right)
    \end{equation*}}
    需 $P_z \gtrsim 2$--$3\GeV$ 以控制 $\lqcd^2/P_z^2$ 幂次修正。

    \item \textbf{短程因子化 / 赝PDF方法 (Radyushkin, 2017)}

    不先做 Fourier 变换，以 Ioffe 时间 $\nu = P_z z$ 为变量，在 $|z| \ll 1/\lqcd$ 下利用 OPE 因子化：
    \begin{equation*}
      \mathfrak{M}_g(\nu, z^2) = \int_0^1 dx\; K(x\nu, z^2\mu^2)\; g(x, \mu^2) + \mathcal{O}(z^2\lqcd^2)
    \end{equation*}
    需 $z \lesssim 0.2\fm$ 以控制 $z^2\lqcd^2$ 高阶扭度修正。
  \end{enumerate}

  \textbf{理论等价性：}在 $\msbar$ 方案下 $C_{gg}$ 与 $K$ 互为 Fourier 变换对，
  $\tilde{g}(x,P_z) = \int \frac{d\nu}{2\pi} e^{ix\nu} \mathfrak{M}_g(\nu, [\nu/P_z]^2)$。
  两种方法交叉验证为格点PDF计算提供内在一致性检验。\vspace{-14pt}
\end{frame}

% ============================================================
% Slide 8: 工作流 I —— 准PDF方法的六步流程
% ============================================================
\begin{frame}{工作流 I —— 准PDF方法 ($\LaMET$) 的完整流程}

  以下给出从格点数据到光锥 PDF 的完整计算链路：

  \begin{description}
    \item[步骤 1] \textbf{构造胶子准算符。} 在格点上定义
    $O_g(z) = F_a^{3\mu}(z)\, \mathcal{W}[z,0]\, F_{\mu}^{a\,3}(0)$，
    其中 $F_{\mu\nu}$ 使用 $\mathcal{O}(a)$ 改进的 clover 离散化。

    \item[步骤 2] \textbf{计算核子三点关联函数。} 在 boost 动量 $\mathbf{P} = (0,0,P_z)$ 的核子态中：
    $C_{3\mathrm{pt}}(\mathbf{P}; t_{\mathrm{sep}}, \tau) = \sum_{\mathbf{x},\mathbf{y}} e^{-i\mathbf{P}\cdot\mathbf{x}}
    \langle N(\mathbf{x}, t_{\mathrm{sep}})\, O_g(\tau;\mathbf{y})\, \bar{N}(\mathbf{0}, 0) \rangle_G$。

    \item[步骤 3] \textbf{提取基态矩阵元。} 利用三点与两点关联函数比值的 $t_{\mathrm{sep}} \to \infty$ 极限：
    $h_g(z, P_z) = \langle P| O_g(z) |P \rangle / (2P_z)$。使用多态拟合法控制激发态污染。

    \item[步骤 4] \textbf{Fourier 变换得准PDF。}
    $\tilde{g}(x, P_z) = \int_{-\infty}^{\infty} \frac{dz}{2\pi}\, e^{i x P_z z}\, h_g(z, P_z)$。
    实际格点上 $z$ 取值范围有限（$|z| \lesssim L/2$），需使用 advanced Fourier 方法处理截断。

    \item[步骤 5] \textbf{非微扰重整化。} Wilson 线引入线性发散 $\sim e^{-\delta m|z|/a}$，
    需在 RI/MOM 方案或比值方案（$h_g(z,P_z)/h_g(z,0)$）下非微扰减除，
    再转换到 $\msbar$ 方案。

    \item[步骤 6] \textbf{匹配到光锥PDF。}
    $g(x, \mu) = \int_x^1 \frac{dy}{|y|}\,
    C_{gg}(x/y, \mu/P_z)\, \tilde{g}_R(y, P_z)$，
    其中 $C_{gg}^{(1)}$ 已知到 NLO (Wang--Zhao 2018; Zhang et al.\ 2019)。
  \end{description}
\end{frame}

% ============================================================
% Slide 9: 工作流 II —— 赝PDF方法与方法对比
% ============================================================
\begin{frame}{工作流 II —— 赝PDF方法及其与 $\LaMET$ 的互补性}

  \textbf{赝PDF方法的四步流程：}

  \textbf{步骤 1.} 用与准PDF相同的矩阵元 $h_g(z, P_z)$，以 Ioffe 时间 $\nu = P_z z$ 重组数据：
  $\mathcal{M}_g(\nu, z^2) = h_g(z, P_z)|_{\nu = P_z z}$。固定 $z$ 下扫描 $P_z$，得 $(\nu, z^2)$ 平面上的二维分布。

  \textbf{步骤 2.} 构造约化 Ioffe-time 分布 (reduced ITD) 以消去主导紫外发散：
  $\mathfrak{M}_g(\nu, z^2) = \frac{\mathcal{M}_g(\nu, z^2)}{\mathcal{M}_g(0, z^2)} \cdot \frac{\mathcal{M}_g(0, 0)}{\mathcal{M}_g(\nu, 0)}$。

  \textbf{步骤 3.} 在 $|z| \lesssim 0.2\fm$ 的短程区域，利用 OPE 因子化：
  $\mathfrak{M}_g(\nu, z^2) = \int_0^1 dx\, K(x\nu, z^2\mu^2)\, g(x,\mu^2) + \mathcal{O}(z^2\lqcd^2)$。

  \textbf{步骤 4.} 选定一组 $z$，用神经网络或参数化 Ansatz 拟合，从积分方程中反演提取 $g(x,\mu^2)$。

  \vspace{6pt}
  \textbf{两种方法的系统学对比：}
  \begin{center}
  \begin{tabular}{p{0.4\textwidth}|p{0.4\textwidth}}
    \hline
    \textbf{准PDF (LaMET)} & \textbf{赝PDF} \\
    \hline
    先 Fourier 变换 $\to$ 后匹配 & 先匹配 $\to$ 后 Fourier 变换 \\
    截断效应来自有限 $z$ & 截断效应来自有限 $\nu$ \\
    幂次修正 $\propto \lqcd^2/P_z^2$ & 幂次修正 $\propto z^2\lqcd^2$ \\
    极限：$P_z \to \infty$（大动量） & 极限：$z \to 0$（短程） \\
    \hline
  \end{tabular}
  \end{center}

  两种方法的\textbf{并行发展}为格点 PDF 计算提供了内在的交叉检验机制——同一物理量由两条独立路径得到一致结果时，可信度大幅提升。
\end{frame}

% ============================================================
% Slide 10: 当前进展、数值挑战与展望
% ============================================================
\begin{frame}{当前进展、数值挑战与展望}

  \textbf{近期格点胶子PDF重要成果：}
  \begin{itemize}
    \item \textbf{LP$^3$ Collaboration (2018--2021)：} 首次格点胶子螺旋度 PDF $\Delta g(x)$。使用 clover-on-clover 算符，验证了在格点上计算胶子分布的可行性。
    \item \textbf{MSULat (2021--2024)：} 首次非极化胶子准PDF 的系统计算，$a \approx 0.09\fm$，$m_\pi \approx 310\MeV$，匹配到 $\msbar$ 后得到唯象上合理的 $g(x)$。
    \item \textbf{HadStruc (2023)：} 蒸馏 (distillation) 方法应用于胶子连通图，信号质量显著提升。
    \item \textbf{$\chi$QCD (2024)：} 混合方案 (hybrid scheme) 重整化应用于胶子准PDF。
  \end{itemize}

  \textbf{走向精确计算面临的四大数值挑战：}
  \begin{enumerate}
    \item \textbf{信噪比指数衰减：} $\mathrm{SNR} \sim e^{-(m_N - \frac{3}{2}m_\pi)t}$，大动量下信号迅速淹没于噪声。
    \item \textbf{激发态污染：} boosted 核子的基态-激发态质量间隔缩小，需 $t_{\mathrm{sep}} \gtrsim 1.0\fm$。
    \item \textbf{多重标度分离：} 需同时满足 $a \ll z \ll 1/\lqcd$ 且 $aP_z \ll 1$，对格距要求极高（$a \lesssim 0.06\fm$）。
    \item \textbf{大动量有限体积效应：} 需 $P_z L \gtrsim 6$--$8$，大体积高动量模拟的计算成本急剧增长。
  \end{enumerate}

  \textbf{展望：} 在物理 $\pi$ 质量、精细格距、大动量下进行胶子 PDF 的系统计算是当前国际竞争的焦点。
  这些结果为即将到来的电子-离子对撞机 (EIC) 提供第一性原理的理论输入，最终目标是系统误差控制在 $10\%$ 以内。

  \begingroup\scriptsize
  \textbf{核心参考文献}\par\vspace{1pt}
  [1] Collins, Soper, Sterman, \textit{Adv. Ser. Direct. High Energy Phys.} \textbf{5}, 1 (1989).
  \hspace{1.5em} [2] X. Ji, \textit{Phys. Rev. Lett.} \textbf{110}, 262002 (2013).\par
  [3] A. Radyushkin, \textit{Phys. Rev. D} \textbf{96}, 034025 (2017).
  \hspace{1.5em} [4] W. Wang, S. Zhao, \textit{Phys. Rev. D} \textbf{97}, 074027 (2018).\par
  [5] Y.-S. Liu et al.\ (LP$^3$), \textit{Phys. Rev. Lett.} \textbf{125}, 232003 (2020).
  \hspace{1.5em} [6] Z. Fan et al.\ (MSULat), \textit{Phys. Rev. D} \textbf{104}, 074507 (2021).\par
  [7] R. Zhang et al., \textit{Phys. Rev. D} \textbf{104}, 034511 (2021).
  \hspace{1.5em} [8] Moch, Vermaseren, Vogt, \textit{Nucl. Phys. B} \textbf{688}, 101 (2004).\par
  \endgroup\vspace{-12pt}
\end{frame}

\end{document}
